\chapter{Newton's Metaphor: Change as Motion \textcolor{red}{(to be
    deleted)}}
\label{cha:newt-metaph-motion}
\markboth{Galileo, Falling Balls, and Newton's Metaphor}{}
% \baselineskip=1.5cm
\section*{Introduction}
% \begin{wrapfigure}[]{r}{2in}
%   \captionsetup{labelformat=empty}
%   \centerline{\includegraphics*[height=2.1in,width=2in]{../Figures/Galileo}}
%   \label{fig:GalileoPortrait}
%   \caption{\textcolor{blue}{Galileo Galilei (1564-1642)}}
% \end{wrapfigure}
Although Newton originally invented and used Calculus to describe the motion of
objects, it was quickly understood that it is a far more general
tool. But dynamics (motion) still provides the essential metaphor for
thinking about phenomena that change in time. We will talk more this
at the end of this chapter.

For now we will take a close look at the original application:
dynamics.  If physics isn't your particular interest be assured that
this is only a starting point.

% It is a truth universally acknowledged that the Calculus was the
% invention (discovery?) of two men, Isaac Newton $(1643-1727)$ and
% Gottfried Wilhelm von Leibniz $(1646-1716).$ Newton developed his
% ideas at his home of Woolsthorpe Manor in Lincolnshire while Cambridge
% University was closed by a plague in the summer of $1665.$ In a period
% of less than two years he made the extraordinary advances while still
% under the age of $25.$ Independently Leibniz developed his ideas
% during the period 1673-1675 while serving as a diplomat in Paris. He
% wrote the first paper\footnote{The rather cumbersome title of Leibniz
% $1684$ paper: \emph{A New Method for Maxima and Minima, as Well as
% Tangents, Which is Impeded Neither by Fractional nor Irrational
% Quantities, and a Remarkable Type of Calculus for This,} is very
% revealing.  Clearly, the problem of optimizing (finding maxima and
% minima) is of central importance and somehow finding tangent lines
% to curves is involved.  That this is a \underline{\emph{new}} method
% means that other methods were available first.} about calculus in
% $1684.$



% These are interesting on their own. However Newton observed that many
% physical problems that are not about velocity or acceleration can be
% analyzed as if they are. That is, the mathematics is the same. We need
% only change the way we think about the meaning of our symbols.



\section{The Simplest Case: A Dropped Ball}
\label{sec:simpl-case:-dropp}
\markright{A Dropped Ball}

To begin we will only consider the case of an object falling
vertically under the influence of gravity.  Suppose we drop a ball
from a height of $9.8$ meters. Will it hit the ground before the end
of that first second?


% \begin{mynotation}{Function}
%   Since $p$ and $t$ are the quantities of interest, and $p$ depends on
%   $t$ we will use the functional notation, $p(t)$ to represent all of
%   this succinctly. Thus $\left.p(t)\right|_{t=0}=p(0).$ This is just a
%   short way to indicate that position $p,$ depends on the elapsed time
%   $t.$
% \end{mynotation}


If we now allow time to tick forward by an infinitesimal amount,
$\dx{t}$, there will be a corresponding infinitesimal change, $\dx{p},$
in $p$. Moreover if $\dx{p}$ is measured in meters and $\dx{t}$ is
measured in seconds then $\dfdx{p}{t}$ is measured in meters per
second, or $m/s.$ That is, the derivative of the ball's position is
its velocity\footnote{We're being a little overly dramatic here. We've
  already seen that the derivative measures the rate of change of one
  quantity with respect to another. This is just a special case of
  that observation.}! Symbolically,
$$
\dfdx{p}{t}=v.
$$
But we also know that $v=9.8t$ so
$$
\dfdx{p}{t}=9.8t.
$$

Hey! Look at that! We have another differential equation to solve. But
unlike some of the differential equations we've seen before this one
is pretty easy to solve. If $\dfdx{p}{t}=9.8t$ then we need only run
the Power Rule backwards to get $p(t) = 9.8\frac{t^2}{2}=4.9t^2$.


So we can now solve our problem, right? All we have to do is find the
time, $t,$ when our position $p(t),$ is zero. That is, we solve
\begin{align*}
  p(t)&=0\\
  4.9t^2&=0\\
  t^2&=0\\
  t&=0.
\end{align*}
So the ball reaches the ground after zero seconds.

\alert{Wait!} That can't be right.

At zero seconds we've just released the ball $9.8$ meters above the
ground so our position at $t=0,$ ($p(0)$), should be $9.8$ not
zero. And $p(t)$ should be zero for some \emph{positive\footnote{Zero
    is neither positive nor negative.}} value of $t,$ since the ball
strikes the ground some time \emph{after} we drop it.  Something seems
to have gone terribly wrong with our analysis of this problem. What
could it be\footnote{Most likely you already see what's gone wrong. If
  so, good for you! If not, no big deal. Keep reading.}?

There really isn't anything wrong with our \emph{reasoning.} The
difficulty is with some of the \emph{assumptions} we made before we
started.  Specifically, when we decided that the ball's position is
given by $p(t)=4.9t^2$ we implicitly assumed that this will give us
the ball's distance from the ground. But it doesn't. We need to
remember that position is always given as a distance \emph{from some
  fixed location}. In this case our $p(t)$ gives us the ball's
distance from its starting position, which is $9.8$ meters above the
ground. Thus when we solved $p(t)=0$ what we found was how long it
takes for the ball to reach its starting position. Since it is already
there, naturally this is $t=0$ seconds.

So we correctly solved the problem that we \emph{stated:} When is
$p(t)=0$? But our implicit assumption led us to state the \emph{wrong
  problem}.

Apparently what we need to solve is $p(t)=9.8.$ This will tell us how
long it takes the ball to travel the $9.8$ meters to the
ground. Solving this gives
\begin{align*}
  p(t) &= 9.8\\
  4.9t^2&=9.8\\
  t^2&=2\\
  t&=\sqrt{2}\approx 1.414
\end{align*}
or slightly less than one and a half seconds.

\begin{Digression}[Coordinate Systems]
%  \noindent\underline{\bf Coordinate Systems}\\
  It is tempting to leave things there, but this problem of implicit
  assumptions -- assumptions we are not aware we are making -- is a
  difficulty that will come back over and over again, sometimes in
  much more complex situations, so we will take a few pages to
  understand it now in this much simpler context.


  When we observed that the velocity $v(t),$ is the derivative of
  position, $p(t),$ we asked what $9.8t$ is the derivative of. The
  answer we \emph{found} was $p(t) =4.9t^2.$ But this is only one
  possibility. The derivative of $4.9t^2+5$ is also $9.8t;$ and of
  $4.9t^2+12;$ and of $4.9t^2-700.$ In fact if $k$ is any constant
  whatsoever then the derivative of $p(t)=4.9t^2+k$ is still $9.8t.$
  If $p(t)=4.9t^2+k$ then the position of the ball ($p(t)$) starts at
  $p(0)=k,$ which is $k$ units away from the origin (the reference
  point).

  When we chose $p(t)=4.9t^2$ as our position function we implicitly
  assumed that $k=0.$ However we could have chosen $k$ to be any
  constant whatsoever. All that changes is what distance is being
  measured by $p(t).$

  % \begin{wrapfigure}[]{r}{1in}
  %   \captionsetup{labelformat=empty}
  %   \includegraphics*[height=2in,width=1in]{Figures-HandDrawn/FallingBall}
  %   \label{fig:FallingBall1}
  % \end{wrapfigure}
  For example either $p_1(t) =4.9t^2$ or $p_2(t)=4.9t^2-9.8$ will work
  as a position function. What is the difference between them?

  If we use $p_1(t) = 4.9t^2$ we are assuming that at time $t=0$ our
  distance from our reference point (origin), $p_1(0)$, is also
  zero. That is, we are starting at the origin. If we use
  $p_2(t)=4.9t^2-9.8$ we are assuming that at time $t=0$ our distance
  from our reference point (origin) is at $p_2(0)=-9.8$ meters. That
  is, our starting position is $9.8$ meters from the origin in the
  negative direction.

  Physically, nothing has changed. But our
  coordinate system is different in  each situation as seen here:\\
  \centerline{\includegraphics*[height=2.5in,width=4in]{../Figures/FallingBall1-2}}
  In the diagram above both images represent the same falling ball.
  At the left $p_1(0) =0$ and $p_1(\sqrt{2})=9.8.$ At the right
  $p_2(0) =-9.8$ and $p_2(\sqrt{2})=0$. The numbers are different, not
  because the physical situation has changed, but simply because we've
  changed the point of reference for the problem.

  In the situation at the left we've implicitly assumed that at time
  zero our starting position is zero. After releasing the ball time
  moves forward as usual so $t$ increases through positive numbers,
  and since $p_1(t)$ measures the distance from our starting point, it
  also increases through positive numbers.

  On the other hand, if we use $p_2(t)=4.9t^2-9.8$ to measure distance
  then we have implicitly assumed that our starting position $p_2(0)$
  is $9.8$ meters above the ground. Notice that the number, $9.8$,
  tells us the distance we are from the origin, but the sign, negative
  in this case, tells us the direction. Thus $p_2(0)=-9.8$ tells us
  that when we release the ball (at time zero) it is 9.8 meters from
  the ground in the \alert{negative} direction. We don't normally
  associate negative numbers with the ``up'' direction, but this is
  just a convention. Nature doesn't care whether we call ``up''
  negative, positive, or Tuesday.

\begin{ProficiencyDrill}{}
  If other values of $k$ are used then the coordinates of the starting
  position (position when $t=0$) and ending position (position when
  $t=\sqrt{2}$) of the ball will change, but the falling time will
  not. Find the position function for each of the following values of
  $k.$ Also give a sketch like the two above for each value of $k,$
  and show that in every case the falling time is $\sqrt{2}.$
  \begin{enumerate}[label={\bf{}(\alph*)}]
  \begin{multicols}{2}
    \item $k=5$
    \item $k=12$
    \item $k=-700$
    \item $k=9.8$
    \item $k=-9.8$
  \end{multicols}
\end{enumerate}
Which of these seems to you like the best choice\footnote{There is
    no correct answer to this question. That is the point.}.
\end{ProficiencyDrill}


% But just because this is a common habit of mind does not mean that
% it is necessarily \emph{always} the best way to measure position.
% For example, one consequence of this for our current problem is that
% the ball falls in the \emph{positive} direction. This is true
% whether we use $p_1(t)$ or $p_2(t)$ In the former case is that as
% $t$ increases the position, $p_1(t),$ changes from zero to $9.8$
% meters -- it becomes ``more positive.'' In the latter it also
% becomes more positive but in this case $p_2(t)$ changes from $-9.8$
% to zero meters\footnote{Perhaps we should call this ``less
% negative?''}. In either case the position now is represented by a
% number that is greater than the position previously, so the positive
% direction is \emph{downward} when we use either $p_1(t)$ or
% $p_2(t).$

\begin{ProficiencyDrill}{}
  It probably feels more natural to use $p(t)=4.9t^2.$ This is because
  we habitually take our current position as the reference point
  (origin). We do this every time we tell a friend something like,
  ``You're about five miles away from me.''

  But we don't normally think of down as positive. Find a new position
  function $p_3(t)$ which models the same falling ball but has down as
  the negative direction. There is more than one way to do this so
  verify that you have a valid function by solving for the time needed
  to reach the ground. You should still get $\sqrt{2}$, just like
  before.
\end{ProficiencyDrill}

Ultimately this all comes down to choosing a coordinate system to work
in. You are probably used to thinking of coordinates as an ordered
pair of numbers $(x,y),$ that specify a location in the plane but the
fact is that you were using a coordinate system long before you
started graphing equations in the plane.  Choosing a coordinate system
consists, essentially, of choosing a distinguished point as a
reference from which the location of all other points is measured. So
as soon as you were taught to think of the real numbers ``as points
strung out on a line'' you were using a coordinate system. The point
at zero is specified (and called the origin) and all other points are
named according to how far, and in what direction, each is from the
origin.

This is what we did when we chose the {\bf k}onstant above. When we
chose $k=0$ so that $p(t)=4.9t^2$ we placed the origin at $9.8$ meters
above the surface of the earth. When we chose $k=-9.8$ so that
$p(t)=4.9t^2-9.8$ we placed the origin at the surface of the
earth. But of course neither gravity nor the ball care what coordinate
system we use so it falls to the ground in $\sqrt{2}$ seconds no
matter where we decide to place the origin.

In the first case since $p(t)=4.9t^2$ we see that $p(0) = 0$ and
$p(\sqrt{2}) = 9.8$ so the ball moved from $0$ to $9.8$ meters. It
follows that ``down'' must be the positive direction and ``up'' the
negative direction.

Similarly if we use $p(t)=4.9t^2-9.8$ we see that $p(0) = -9.8$ and
$p(\sqrt{2}) = 0.$ In this case, although the position of the ball is
always negative, it moves from a more negative position, $-9.8,$ to a
non-negative position, $0.$ So again ``down'' must be the positive and
``up'' the negative direction.

This seems natural when we think about a falling ball, but of course
when we choose a coordinate system we get to choose both the origin
and the directionality. In the next section it will be more natural to
take ``up'' as the positive direction.
\end{Digression}


\section{Falling Balls: How High Will the Ball Go?}
\label{sec:tossed-ball}
\markright{How High Will The Ball Go?}
Instead of simply dropping the ball suppose we toss it vertically
upwards with some initial velocity of say \(15 m/s.\) Then after one
second it will have \emph{lost} \(9.8 m/s\) of its initial
velocity. That is, its velocity after \(1\) second will be
\(15-9.8=5.02\,m/s.\) After \(2\) seconds, its velocity will be
\(15-9.8\times{}2=-4.6 m/s.\)

In general, after \(t\) seconds its velocity will be \(15-9.8t\, m/s.\) 
Notice that the velocity is positive at the one second mark $(t=1)$
and negative at the two second mark $(t=2.)$

As we observed at the end of the last section, it seems natural to
interpret \underline{positive} velocity as climbing (moving in the
positive direction). We'll need to keep this in mind as we proceed.

So the height of the ball increases until its velocity drops to
zero. At that time the ball's height has reached its zenith. After
that the ball will drop. But now its velocity will become larger in
the negative direction (because the ball is falling -- moving in the
negative direction), Although we don't know exactly what the height is
at any given time we can see that there is a time, somewhere between
one and two seconds after the ball is released when it reaches its
zenith. At that time the ball's velocity will have decreased to zero
because for an instant at the top of its climb the ball is neither
rising nor falling. Its velocity is
zero. Setting the velocity equal to zero and solving for $t$ gives:
\begin{align*}
  v(t)&=0\\
  15-9.8 t&=0 \\
  t=15/9.8&\approx 1.53 
\end{align*}
Thus the  maximum height is reached after approximately $1.53$ seconds.

\begin{embeddedproblem}{}
  Suppose the initial velocity is double what it was originally.
  Would it take twice as long for the ball to reach its maximum
  height?  What if the initial velocity was half?
\end{embeddedproblem}

So far, so good.  But now we ask: How high will the ball go? 

Once again what we need is not the velocity function, $v(t) =15-9.8t,$
but the position function, $p(t).$ Once we have that, since we know
how long it takes to reach the maximum height ($\approx 1.53$
seconds), all we need to do is compute $p(1.53).$

Since we know that position is the derivative of velocity all we need
to do is ask ourselves; What is $v(t) = 15-9.8t$ the derivative of?

\begin{ProficiencyDrill}
  Give this a few moment's thought before reading on. See if you can
  figure it out for yourself. We'll wait. \hint{Run the Power Rule
    backwards.}
\end{ProficiencyDrill}

It should be clear that $p(t) = 15t-4.9t^2$ could  be the  position
function. 
Is it?


Wait a moment! Remember the trouble we had before with the
{\bf{}k}onstant, $k?$ Can't the same thing happen here? Let's think
about this a bit more.

As before the derivative of $ p(t) = 15t-4.9t^2 +k $ is
$ v(t) =15-9.8t $ no matter what the numerical value of $k$ is, so
let's see if we can find the value of $k$ that makes things easiest
for us. If\footnote{Notice that we've rearranged the order of the
  terms. This isn't really necessary, but it's a good idea to always
  write polynomial in the standard form, unless there is a compelling
  reason not to. This makes things easier to read when you come back
  looking for the errors that will inevitably happen.}
$p(t) = -4.9t^2+15t +k$ then at time $t=0$ the ball's position is
$p(0) = -4.9(0)^2+15(0) +k = k.$ Since we're standing on the ground
when we toss our ball we'll probably release the ball about $7$ feet
above the ground\footnote{We (the authors) are both a little over six
  feet tall and we'd release the ball a little above our heads. For
  the sake of keeping things simple we'll say this is $2$ meters. If
  you are much taller or much shorter you should adjust this number as
  appropriate.}, which we'll round up to  $2$ meters. If we set $k=2$ then we've
placed the origin at ground level since $p(0)=2$. In words, we start
out $2$ meters from our reference point, the ground. Thus our
position function is
$$
p(t) = -4.9t^2+15t +  2
$$
and therefore the maximum height of the ball is approximately $p(1.53) =
-4.9(1.53)^2+15(1.53) +  2 \approx 13.5$ meters.

\begin{embeddedproblem}{}
  Suppose you toss the ball twice as hard so that its initial velocity
  (its velocity when it leaves your hand) is $30 m/s.$ 
  \begin{enumerate}[label={\bf{}(\alph*)}]{}
  \item How long will it take to reach its maximum height?
  \item What is its maximum height?
  \item How long does it take the ball to hit the ground?
  \item What is its velocity when it hits the ground?
  \end{enumerate}
\end{embeddedproblem}

We now have everything we need to describe the motion of a ball thrown
vertically. We have:
\begin{enumerate}
\item The derivative of the position $p(t),$ is the velocity $v(t).$
\item The derivative of the velocity $v(t),$ is the acceleration.
\item At the surface of the earth the acceleration due to gravity is
  $9.8\,m/s^2.$ To keep our analysis as general as possible we will
  just say that the acceleration is a constant and use $g$ to
  represent it, as is usual.
\end{enumerate}

First, to find the velocity function $v(t),$ we ask, what is the
constant $g$ is the derivative of?  Clearly this is $v(t) =g t+ v_0.$
To determine the value of $v_0$ we evaluate the velocity when $t=0.$
Since
$$
v(0) = g\cdot0+v_0 = v_0
$$
we see that $v_0$ is the \emph{initial velocity} of the ball.

Do you recognize what we just did? We solved another Initial Value
Problem (IVP). The differential equation was:
$$
\dfdx{v}{t}=g
$$
and the Initial Condition was $v(0)= v_0.$ It is not a particularly
difficulty IVP, but it is worth noting that these things keep coming
up.

Next we find the position function, $p(t)$, by asking what 
$$
v(t) = g t+v_0
$$ 
is the derivative of? Clearly this is 
$$
p(t) = 
\left(
  \frac{g}{2}
\right)t^2+v_0t+p_0
$$
where $p_0$ is the initial \emph{position} by a similar argument.

Once again this is an IVP. In this case it is the IVP
$$
\dfdx{p}{t}=gt+v_0,\ \ p(0)=p_0.
$$

\begin{myexample}
  Suppose you toss the ball vertically with an  initial velocity of
  $30\,m/s$ while standing at the edge of a $120$ meter building
  is $30 m/s.$ 
  \begin{enumerate}[label={\bf{}(\alph*)}]
  \item How long will it take to reach its maximum height?
  \item What is its maximum height?
  \item How long does it take the ball to hit the ground?
  \item What is its velocity when it hits the ground?
  \end{enumerate}
\end{myexample}

\subsection{Lateral Motion, and Parametric Functions}
\label{sec:lateral-motion}
\begin{wrapfigure}[]{r}{2in}
  \captionsetup{labelformat=empty}
  \includegraphics*[height=2in,width=2in]{../Figures/FallingBall1}
  \label{fig:FallingBall}
\end{wrapfigure}
As we saw if a ball is tossed vertically from a height of $2$ meters above the
ground then its \emph{vertical position} is given by:
$$
p(t) = -4.9t^2+15t +  2
$$
and the  graph of this function is shown at the right. 

Before we go on though, let's take a moment to notice that 
although this graph contains a great deal of information about
the ball's motion it is \underline{not} a picture of the ball's
\emph{actual path} through the air. If it is thrown vertically then
the ball moves vertically upward from our hands and then vertically
downward to the ground. Since the horizontal axis of this graph
represents the elapsed time, $t,$ this graph is actually a fairly abstract
representation of the purely vertical motion of the ball.  If we start
our clock at the moment the ball leaves our hands then at time $t=1$
we have $p(1)=12.1$ and the point $(1,12.1)$ represents the fact that
one second after the ball is released its height is $12.1$ meters
above the ground. It does \emph{not} say that the ball is at some
position labeled $(1,12.1).$ This is clear from the labels $p(t)$
(position) and $t$ (time) on the vertical and horizontal axes.


However, we all know that it is nearly impossible to toss a ball
perfectly vertically. Inevitably the ball will receive some non-zero
lateral (horizontal) velocity. When that happens the path the ball
follows will indeed look very much like the previous graph. This can
lead to  confusion so it was worth taking a few moments to clarify this
a bit. 

When the ball moves perfectly vertically the situation is as we've
described, and its \emph{vertical} position is given by
$$
y(t) = -4.9t^2+15t+2,
$$
(where we've renamed the position function to $y(t)$ to emphasize that this
gives the vertical position only). Thus $\dfdx{y}{t} = -9.8t +15$ represents the
vertical velocity of the ball, as before. We saw earlier that the
ball will reach its apex  when $t=1.53$ seconds and this is consistent
with what we see in our graph.

If we want to learn how  long will it take for the ball to reach the
ground we set the vertical coordinate equal to zero and solve for $t$,
$$
y(t) = -4.9t^2+15t+2 = 0.
$$
\begin{ProficiencyDrill}
  Show that $t\approx 3.2.$  
\end{ProficiencyDrill}
And now that we know when the ball will strike the ground we can
compute its vertical velocity at that time as well:
$$\eval{\dfdx{y}{t}}{t}{3.2} = -16.36\text{ m/sec.}$$
But remember that the ball is moving straight up and down. There is no
side to side motion at all, so its \emph{horizontal velocity} is zero
because it is not moving horizontally at all.

\begin{wrapfigure}[]{l}{2in}
  \captionsetup{labelformat=empty}
  \centerline{\includegraphics*[height=1.5in,width=2in]{../Figures/FallingBall2}}
  \label{fig:FallingBall2}
\end{wrapfigure}
Now suppose (more realistically) that we release the ball with an
initial \underline{horizontal} velocity is $5 m/s.$ Then the
\emph{horizontal} position of the ball will be
$$
x(t)=5t.
$$
and the \emph{horizontal} velocity will be by $\dfdx{x}{t}=5$
m/sec. That is, the ball will continue moving horizontally at $5$
m/sec throughout its fall.  If we want to create an actual picture of
the path the ball follows we need to plot its spatial coordinates,
$(x(t), y(t))$ at every moment, $t$, of its flight, as in the sketch
at the left. Notice that the axes are labeled $x(t)$ and $y(t)$, both
of which measure distance. In this graph the horizontal and vertical
coordinates represent the physical, distance the ball has traveled,
horizontally and vertically, and the point $(x(2), y(2)) = (10, 12.4)$
\underline{does} represent the physical location of the ball.

\begin{myexample}
  The path of a weight rotating in a circle with radius $= 1$ in the
  way we just described is described by setting $x(t)$ and $y(t)$ as
  follows:
  \begin{align*}
    x(t)&=-\cos(t)\\
    y(t)&=\sin(t).
  \end{align*}
  When $t=\pi/4$ we have $(x(\pi/4), y(\pi/4) =
  (-\cos(\pi/4),\sin(\pi/4) = (-1/\sqrt{2},1/\sqrt{2})$, so the weight
  is at the base of the red tangent arrow in our sketch. If everything
  we've just said is correct then it should be true that
  $
  \eval{\dfdx{y}{x}}{t}{\pi/4} = 1
  $, the slope of the red tangent line.  Let's see if it is.
  $$
  \eval{\dfdx{y}{x}}{t}{\pi/4}=\eval{\frac{\dfdx{y}{t}}{\dfdx{x}{t}}}{t}{\pi/4}=\eval{\frac{\cos(t)}{\sin(t)}}{t}{\pi/4}=\cot(\pi/4)=1.
  $$
\end{myexample}
\begin{embeddedproblem}{}
  For each set of conditions given below, answer the following eight questions:
  \begin{description}
  \item[{\bf Question 1:}] What are the $x$ and $y$ coordinates of the ball's
    position at time, $t.$
  \item[{\bf Question 2:}] Plot the ball's position on the $x-y$ plane.
  \item[{\bf Question 3:}] Plot the \underline{vertical position only} of the ball
    on the $p-t$ plane. This will not necessarily look the same
    as the graph you drew for the previous question\aside{You should be curious about
    the difference. In fact, you should be so curious that you will
    spend a few minutes trying to puzzle out why and how they are
    different. If you are not curious, pretend that you are and do
    it anyway. Learning happens when you try, not when you succeed.}. 
    \item[{\bf Question 4:}] At what time does the ball reach its maximum height?
    \item[{\bf Question 5:}] What is the maximum height?
    \item[{\bf Question 6:}] How long before the ball hits the ground?
    \item[{\bf Question 7:}] Where does it strike the ground?
    \item[{\bf Question 8:}] What is it's horizontal and vertical velocity at the moment
      it strikes the ground?
    \end{description}
    
  \begin{enumerate}[label={\bf{}(\alph*)}]
  \item Suppose our initial horizontal position is $2 m/s,$ and that
    the initial vertical velocity is $15 m/s.$
  \item Suppose our initial vertical velocity is $\frac{1}{2} m/s,$ and that
    the initial vertical velocity is $15 m/s.$
  \item Suppose our initial horizontal velocity is $0 m/s,$ that
    the initial vertical velocity is $15 m/s,$ but that the ball is
    accelerating to the right at a constant rate of $2 m/s^2.$
  \end{enumerate}
\end{embeddedproblem}

\subsection*{Parametric Functions}
When the representation of some phenomenon is given as two numbers,
($x$ and $y$ in this case) which each depend on a third ($t$ in this
case) the underlying variable is called a parameter and the
representation is known as \term{parametric}. Thus the following is the
parametric representation of a unit circle:
\begin{align*}
  x(t) &= -\cos(t)\\
  y(t) &= \sin(t).
\end{align*}
In this representation we have implicitly assumed that time ($t$) will
continue indefinitely so we keep going round and round forever.

Of course that can't happen. Eventually we'll get tired and stop
swinging the weight around and around. If we stop after two go-rounds
then we need to add the constraint
$$
0\le t\le 4\pi.
$$

\begin{myexample}
  So to represent a single unit circle starting at the point $(-1,0)$
  and proceeding clockwise we have the parametric \term{function}
  (representation)
  \begin{align*}
    x(t) &= -\cos(t)\\
    y(t) &= \sin(t)\\
    0 \le & t \le 2\pi.
  \end{align*}
  The first two formulas are the \term{parametric equations}. Taken
  together they form a single function:
  $$
  \phi(t)=
  \begin{pmatrix}
    -\cos(t)\\
    \sin(t)
  \end{pmatrix}
  $$
  The third formula, $0 \le  t \le 2\pi$, defines the
  \term{domain} of the function, $\phi$.
\end{myexample}
We will formalize all of this with the following  definition.\\
\begin{mydefinition}[Parametric Functions]
  \label{def:parametric-functions}
  Given a \term{parameter} $t$, where $a\le t \le b$, the expression
  $$
  \phi(t) =
  \begin{pmatrix}
    x(t)\\
    y(t)
  \end{pmatrix}
  $$
  defines a function $\phi:\RR\rightarrow \RR^2.$
\end{mydefinition}

There is a lot that is new in
definition~\ref{def:parametric-functions}, both conceptually and 
in the notation so we will take a moment to discuss both.

First the notation. There are three questions we need to address:
\begin{mynotation}{Column vs. Row}
  \begin{description}
  \item[Why did we write the column:
    $\begin{pmatrix}
      x(t)\\
      y(t)
    \end{pmatrix}$, rather than the row: $(x(t), y(t))?$]
    \ \\
    Honestly, there really is no compelling \emph{mathematical} reason
    to prefer one over the other in this case. Moreover, since you are
    already comfortable with interpreting an ``ordered pair''
    $(x, y)$ as a point on the $x-y$ plane the row for a good case can
    be made for sticking with it.

    Why would we change this on you? Do we simply enjoy making things
    hard?

    No, of course not.
    
    But if we want to teach you to communicate effectively (we do)
    then we need to teach you to use the the notation that has become
    standard and which you will see most often in later courses. For
    better or worse, this is the column notation. It is better for
    you get comfortable with it now while things are still fairly
    simple.
  \item[What does $\phi:\RR\rightarrow\RR^2$ mean?]
    \ \\
    As above, this is a standard notation that you will have to learn
    to use in more complicated situations. Best to get comfortable
    with it now.
    
    Because this is brand new notation for you it probably appears
    complicated, but it really isn't.  This is just the statement that
    our function, $\phi$ picks up one real number ($\RR^1$, or just
    $\RR$) and returns  a point in the plane ($\RR^2$). Of course
    we need two numbers, $x$ and $y$, to specify a point in $\RR^2$.

    We write it in this very compact form for the simple reason that
    eventually you will need to become comfortable with functions that
    return $10$, $1000$, or even $10^{200}$ real numbers. Rather than
    writing out a column of length $10$, $1000$, or $10^{200}$ it is
    much simpler to write
    \begin{align*}
      \phi&:\RR\rightarrow\RR^{10},\\
      \phi&:\RR\rightarrow\RR^{100},\text{ or}\\
      \phi&:\RR\rightarrow\RR^{10^{200}}.
    \end{align*}
    The completely general situation is:
    $$
    \phi:\RR^n\rightarrow\RR^m, n,m \in \NN
    $$
    which just says that a function $\phi$ picks up $n$ real numbers
    (the coordinates of a single point in $\RR^n$) and returns $m$
    real numbers (the coordinates of a single point in $\RR^n$).  The
    good news for you is that in this course $n$ will always be equal
    to $1$ and $m$ will be usually be $1, 2,$ or $3$ (though mostly
    just $1$).
  \item[If $\phi:\RR\rightarrow\RR^2$ returns \emph{two} numbers is it
    \emph{really} a function?]
    \ \\
    Yes, it really is. The reason that this probably feels wrong to
    you is that so far you've mostly been dealing with functions
    defined by something like $f(x)=x^2$. That is,
    $f:\RR\rightarrow\RR, f(x)=x^2$.

    If all you've ever dealt with are functions from $\RR$ to $\RR$
    then it is easy to get the impression that the function and the
    formula are the same thing. They are not. Functions are very
    general objects.

    Fundamentally a function is a way of expressing associations
    between ``things.'' For example the function $f(x)=x^2$
    associates $1/2$ with $1/4$,  $2$ with $4$, $1/3$ with $1/9$ and
    so on. \underline{But} a function must associate one object in
    its domain (stuff coming in) with precisely one object in its
    range (stuff coming out). For example if $2$ is associated with
    both $3$ and $4$ then we don't have a function.

    But wait! Doesn't
    $$
    \phi(t)=
    \begin{pmatrix}
      t+1\\
      t+2
    \end{pmatrix}
    $$
    do  exactly that? After all $\phi(2) =
    \begin{pmatrix}
      3\\4
    \end{pmatrix},$ right? So this is \alert{not} a function?

    Wrong. This is a function.

    What's being associated with $t=2$ is not the two numbers $x=2$
    and $y=4$ but the (single and unique) point in the $x-y$ plane
    with coordinates $\begin{pmatrix}3\\4\end{pmatrix}$. This is a single ``thing'',
    but because it lies in the plane it requires two numbers to 
    identify it.
  \end{description}
\end{mynotation}
% \ \newpage{}
% \ \newpage{}
% \ \newpage{}
% \ \newpage{}
% \ \newpage{}
% \ \newpage{}
% \ \newpage{}
% \ \newpage{}
% To be sure, in later courses you will be dealing with parametric
% representations  with
% more than two coordinate functions like for instance, the function
% $$
% \phi(t)=
% \begin{pmatrix}
%   x(t)\\
%   y(t)\\
%   z(t)\\
%   w(t)\\
%   u(t)\\
%   v(t)\\
% \end{pmatrix}.
% $$
% which is probably easier to read the column version than the row
% version:
% $$\phi(t) = \left(x(t), y(t), z(t), w(t), u(t), v(t)\right),$$ but
% really, we just had to choose one.


% \ \newpage{}
% \ \newpage{}
% \ \newpage{}
% \ \newpage{}
% \ \newpage{}
% \ \newpage{}
% \ \newpage{}
% \ \newpage{}




% \begin{wrapfigure}[]{l}{2in}
%   \captionsetup{labelformat=empty}
%   \centerline{\includegraphics*[height=1in,width=2in]{../Figures/FallingBallPath1}}
%   \label{fig:FallingBallPath1}
% \end{wrapfigure}
% \noindent{}\underline{except} for the labelling on the axes. Notice that the
% horizontal axis shown is labelled ``$x$'' and the vertical is
% ``$y.$'' This is because the horizontal axis now represents the

% The vertical axis still represents the vertical position of the ball
% but since the position now requires two coordinates, $x$ and $y,$ we
% simply label the vertical axis ``$y$'' as per tradition. 

% The reason the two graphs are identical in this case (except for the axis labels,
% of course) is that we have assumed an initial horizontal
% velocity of $v(0)=1 \text{m/sec}$  which corresponds to time moving
% forward at a constant rate of $1\text{m/sec}.$ If we also choose our horizontal origin to be
% where we are standing when we release the ball, then the horizontal
% position is given by:
% $
% x=t.
% $
% Since the initial vertical velocity is unchanged (as is the
% acceleraton due to gravity: $9.8 m/s^2$) the vertical position is
% still 
% $$
% y(t) = 2 + 15t-4.9t^2.
% $$

% We use $y$ (instead of $p$) to represent the vertical position of the
% ball because the ball's position in this situation requires
% two\footnote{Actually we need $3.$ Do you see why we can ignore the
% third this time? If not, pause here and think about it for a moment.
% We'll wait.}
% coordinates, $x$ \underline{and} $y,$ whereas previously we only
% needed one, the position, or $p$, coordinate.

% If we want to use $p$ to represent the ball's position at time $t$ (we
% do) then $p(t)$ will have to return two spatial ($x$ and $y$)
% coordinates for each value of $t:$
% $$
% p(t) = \left(x(t), y(t)\right).
% $$

% To see why the two graphs above coincide we need only observe that
% since 
% $$
% x=t
% $$
% it doesn't matter whether we label the horizontal axis $t,$ in which
% case the graph gives the vertical position at any time $t,$ or $x,$
% in which case the curve gives the actual position of the ball. They
% \underline{look}  the same even if they mean different things.

% We got $x=t$ because we assumed an initial horizontal velocity,
% $v_x(0)$ of $1 m/s$ and an initial horizontal position, $x_0$ of
% zero. 

% % \begin{wrapfigure}[]{l}{2in}
% %   \captionsetup{labelformat=empty}
% %   \centerline{\includegraphics*[height=1in,width=2in]{../Figures/FallingBallPath2}}
% %   \label{fig:FallingBallPath2}
% % \end{wrapfigure}
% Suppose we set our horizontal origin\footnote{\textcolor{red}{The
% figure needs to be adjusted.}} $3$ meters to the left of where
% we released the ball. Then we get
% \begin{align*}
    %     x&=t+3\\
    %     y&=2 + 15t-4.9t^2,
             %   \end{align*}
             %              so the ball follows the given trajectory in this sketch.\\
    %     \centerline{\includegraphics*[height=1in,width=2in]{../Figures/FallingBallPath2}}

The position of a moving object is always a function of time because
that's what position depends on, time.  A location and the time of
arrival at that location are the two things the position function
associates.

But what is its velocity? 
\begin{myexample}
    %     \begin{wrapfigure}[]{r}{2in}
    %     \captionsetup{labelformat=empty}
    %     \centerline{\includegraphics*[height=1.5in,width=2in]{../Figures/Parametric1}}
    %     \label{fig:Parametric1}
    %     \end{wrapfigure}
  Suppose a point is moving in the $x-y$ plane along the path
  $$
  \phi(t) = \begin{pmatrix}x(t)\\y(t)\end{pmatrix}
  =
  \begin{pmatrix}
    10t-3\\     
    2 + 15t-4.9t^2,
  \end{pmatrix}, 0\le t \le 4.
  $$
  The graph\aside{Your teacher may express some annoyance with us
    (the authors) for calling this a ``graph''. And your teacher is
    correct to be annoyed. Properly speaking this should be called
    the ``trace'' or the ``path'' of the function, not its graph. The
    word ``graph'' is reserved (mostly) for plotting the $x-y$
    coordinates of a function: $f:\RR\rightarrow\RR$, just like
    you've always used it. We're using the wrong word here because we
    didn't want to take time to define a new word and concept just to
    talk about this example. Especially when, in our opinion, doing
    so would add nothing to your understanding and probably would
    lead to unnecessary confusion. Still, this is not a graph,
    strictly speaking. Your teacher can explain why not. If you're
    interested, ask them.} of this function looks like the
  black curve in the sketch:\\
  \centerline{\includegraphics*[height=1.5in,width=2in]{../Figures/Parametric1}}
  What is its velocity when $t=2$?

  Actually, this question is too vague to answer.
  The point has several velocities. Since its horizontal
  position is given by $x(t) = 10t-3$ it should be clear that its
  velocity in the horizontal direction is constantly $\dfdx{x}{t}=10.$  Similarly the
  velocity in the vertical direction will be given by
  $\dfdx{y}{t}=\eval{15-9.8t}{t}{2}= -4.6$. As before, since $x(t)$
  and $y(t)$ are positions, their respective derivatives are
  velocities. Since $y(t)$ only gives the vertical position,
  $\dfdx{y}{t}$ only gives the vertical velocity. Similarly for
  $x(t)$.

  But when $t=2$ the point isn't just moving horizontally or
  vertically. It is actually moving in the direction of the tangent
  line at that point, as seen in the sketch. It is traditional to
  denote this by $s$. So the rate of change of $s$ with respect to
  time, $\dfdx{s}{t}$ will be the velocity \emph{in the direction of
    motion.} Sometimes this is called the \emph{tangential
    direction.}  Can we find out how fast it is moving in the
  tangential direction, $\dfdx{s}{t}$?

  Sure. All we have to do is find out what $\dx{s}$ is, and divide
  that by $\dx{t}$.  Consider the differential triangle drawn in
  red. Clearly, by the Pythagorean Theorem,
  \begin{align*}
    \dx{s}^2&=\dx{x}^2+\dx{y}^2.\\
    \intertext{Dividing by $\dx{t}^2$ we have}
    \frac{\dx{s}^2}{\dx{t}^2}&=\frac{\dx{x}^2}{\dx{t}^2}+\frac{\dx{y}^2}{\dx{t}^2}\\
    \left(\dfdx{s}{t}\right)^2&=\left(\dfdx{x}{t}\right)^2+\left(\dfdx{y}{t}\right)^2\\
    \intertext{so that}
    \dfdx{s}{t}&=\sqrt{\left(\dfdx{x}{t}\right)^2+\left(\dfdx{y}{t}\right)^2}
  \end{align*}
\end{myexample}

\begin{myexample}
  We will continue using the same function:
  $$
  \phi(t) =
  \begin{pmatrix}
    10t-3\\     
    2 + 15t-4.9t^2,
  \end{pmatrix}, 0\le t \le 4.
  $$
  but now we will change the way we think about it.  We've been
  thinking of $\phi(t)$ as a function describing the location of a
  point. That is, we've been focused on motion. We can also think of
  the graph above as simply a set of points in the plane. Now nothing
  is moving. We simply have a function that picks up a number between
  $0$ and $4$ and returns the coordinates of the point on the curve
  associated with that number.

  Since there is no value in choosing another variable name we
  continue to use $t$ as our parameter. But keep in mind that
  $\phi(2)$ no longer represents a moment of time, it is now a
  position on a  curve, not a position we are passing through.

  The picture
  remains the same, but with this change in perspective another
  question presents itself: What is the slope of the tangent line at
  $\phi(2)$?
  Clearly it is still ``the infinitesimal change in $y$ over the
  infinitesimal change in $x$'' so it is given by
  $\dfdx{y}{x}$, but how could we actually compute it?

  Actually this is easy! Since $\dx{y}$ appears in the numerator of $\dfdx{y}{t}$ and
  $\dx{x}$ appears in the numerator of $\dfdx{x}{t}$ we see that
  $$
  \frac{\dfdx{y}{t}}{\dfdx{x}{t}} =
  \frac{\dx{y}}{\cancel{\dx{t}}}\cdot\frac{\cancel{\dx{t}}}{\dx{x}}
  =\dfdx{y}{x}.
  $$
  So the slope of the line tangent to $\phi(t)$ at $t=2$ is given by
  $$
  \eval{\dfdx{y}{x}}{t}{2}
  =\eval{\frac{\dfdx{y}{t}}{\dfdx{x}{t}}}{t}{2} =
  \eval{\frac{10}{15-9.8t}}{t}{2} \approx
  -2.17
  $$
  
  It really is that simple, but some care must be taken. Notice in
  particular that although $\dfdx{y}{x}$ is given in terms of $t$ not
  $x$ or $y$. Necessarily so since $t$ is the variable here. Both $x$
  and $y$ are given in terms of $t$ so  the derivative of
  $y$ with respect to $x$ will also be in terms of $t$.
\end{myexample}

\begin{mynotation}{Differential}
  We would be remiss if we did not point out that the more you think
  about the meanings of the preceding equations the less sensible they
  seem to become. In particular, the meaning of $\dfdx{t}{x}$ is very
  unclear. We have come to think of an expression like $\dfdx{y}{x}$
  as representing the rate of change of $y$ with respect to $x$. But
  then, what does it mean to say that $\dfdx{t}{x}$ the rate of change
  of time with respect to the horizontal distance, $x$? The idea seems
  to be meaningless.

  The difficulties here are real, and they are closely tied to the
  problems inherent in the idea of differentials. These issues need to
  be resolved. We will not pretend otherwise. But we do not need to
  resolve them right now so, for now, we will simply ignore them. If
  the expression $\dfdx{t}{x}$ bothers you, good. It should. But for
  now it is safe to think of it in the same way we always have: as a
  ratio of differentials. In this case it is the reciprocal of the
  horizontal velocity, $\dfdx{x}{t}$. So if, for example,
  $\dfdx{x}{t}=5$, then  $\dfdx{t}{x}=\frac{1}{5}$.
\end{mynotation}




\section{Motion as a Metaphor}
\label{sec:newtons-view:-motion}
\markright{Newton's Metaphor}

Newton had a specific purpose in mind when he invented Calculus: He
was building on the work of Galileo and others to describe the
physical motion of the moon and planets under the influence of
gravity. Naturally this affected the way that he thought about
Calculus and the language he used to talk about it. Although he
quickly realized that his Calculus could be used much more widely he
was deeply tied to the physics of kinematics (motion), and the
language that he used reflects this bias.

If we think of the motion as metaphorical rather than literal we can
extend Newton's kinematic approach to problems which do not
necessarily involve motion. To see how this works we will reconsider
example~\ref{example:MaxTriangle} from the Introduction.

\begin{embeddedproblem}{}
  The lengths of two sides of a triangle are $a$ and $b.$ If the third
  side is chosen in such a way that the area of the triangle is as
  large as possible what is the length of the third side?
\end{embeddedproblem}
\begin{wrapfigure}[]{r}{2in}
  \captionsetup{labelformat=empty}
  \centerline{\includegraphics*[height=1in,width=2in]{../Figures/Advice1-1}}
  \label{fig:Advice1-1}
\end{wrapfigure}
If we visualize the triangle as static then we apparently have no
choice but to imagine that we need calculate the areas of all
(infinitely many) possible such triangles and choose the largest one.

A daunting task at best.

On the other hand, since $a$ and $b$ are the legs of a triangle they
\emph{must} share an endpoint. We choose a coordinate system by
designating that shared point to be the origin and then connect the
other two endpoints with a straight line $c$ and it should be clear
that we can generate a triangle by thinking of one leg, say $a,$ as
fixed and the other as rotating about the origin.  In this way we can
introduce motion, at least metaphorically, into what is really a
static problem. Do you see how this helps?

We know that the area of a triangle is
$\frac12\left[\text{$(b)$ase}\times\text{$(h)$eight}\right].$ If we
take $b$ as the base then the height will be $a\sin(\theta)$ where
$\theta$ is the angle between $a$ and $b.$ So for given values of $a$
and $b$ the area depends on -- is a function of -- $\theta.$

And what function is that? Easy!! It is\footnote{
  At this point we can solve the problem without Calculus. Clearly all
  we need to do is find the value of $\theta $ for which $\sin(\theta)$
  is maximal. Everyone knows that $\sin(\pi/2)=1$ and that this is as
  large as the sine function gets. Thus the largest possible area for a
  given $a$ and $b$ is
  $$
  A(\pi/2) = \frac12ab\sin(\pi/2) = \frac12ab.
  $$

  Were you able to intuit the correct answer? Whether you did or not
  take a moment to consider why your intuition took you in the direction
  it did. Why it worked. Or didn't. We'll wait.
}
\begin{align*}
  A(\theta) &=
              \frac12\left[\text{$(b)$ase}\times\text{$(h)$eight}\right]  \\
  A(\theta) &= \frac12ba\sin(\theta).
\end{align*}


Clearly we don't need anything as powerful as Calculus to solve this
problem, but since learning Calculus is what we're here for we'll use
it anyway. Since we now have the area of our triangle given as a
function of the angle $\theta$ we no longer need think of $A$ as the
area, and $\theta$ as an angle. Instead, if we choose to we can think
of $A$ as position and $\theta$ as elapsed time. If we start at time
$\theta=0$ and move forward it is clear that our position, $A,$ will
be as far from the origin as it will get when we stop moving away from
the origin and start moving back towards it. That is, when our
velocity changes from positive to negative. This of course happens
when the velocity is zero. Having introduced the metaphor of motion
into our problem we are now thinking of $A$ as a distance and $\theta$
as time so that $\dfdx{A}{\theta}$ represents velocity just like
before.

And what is our velocity? Why it's the derivative $\dfdx{A}{\theta}$
of course, just like before. Apparently we can find the maximum
distance traveled (maximum area of our triangle) by solving the
equation\footnote{Notice that this is \emph{not} a differential
  equation. This is just a plain, old, algebraic equation of the sort
  your mathematics teachers have been asking you to solve for years. }
$$
\dfdx{A}{\theta} = 0.
$$
Since the only variable is $\theta$ ($a$ and $b$ are constants) $A(\theta)=\frac12ab\sin(\theta)$ we have
\begin{align*}
  \dfdx{
  \left(
  \frac12ab\sin(\theta)
  \right)}{\theta} &= 0 \\
  \frac12ab\cos(\theta) &= 0 \\
  \cos(\theta) &= 0, \\
\end{align*}
and everyone knows that $\cos(\theta)=0$ when $\theta=\pi/2.$ So the
area of the triangle is a maximum when $\theta=\frac{\pi}{2}$ and the
maximum area is $A(\pi/2)=\frac12ab\sin(\pi/2)=\frac12ab.$

Newton and his followers were quite explicit in using this
metaphor. They would have interpreted the formula
$A=\frac12ab\sin(\theta)$ as giving the length of a straight line as
it is traced out by a moving point. The length of the line grows (or
shrinks) in accordance with, the way the area of our triangle $\theta$
grows (or shrinks). Thus they explicitly identified the angle with the
passage of time, and the rate of change of the area of the triangle
with the velocity of the moving point.

In practice we won't need to be quite so explicit about bringing
motion into the problem.  As we say in the title of this section,
motion is really just a metaphor. When we interpret the expression,
$\dfdx{A}{\theta}$ as ``the rate of change of $A$ with respect to
$\theta$'' there is a clear parallel between the way $A$ changes
depending on the angle $\theta$ and the way position changes depending
on elapsed time.


What the first derivative tells us about the motion of an object is
contained in the following theorem. Notice that this theorem is stated
in terms of motion as per the Newtonian viewpoint\notetoself{Starting
  here I've got
  several different statements of the First Derivative Test. This
  seemed like a good idea when I was writing it but now I'm not so
  sure. Should we call all of the early versions something else and
  only all the last one~\ref{thm:FDTproved} the First Derivative Test?
  Note that the last one is slightly different from the others. Or maybe
  we should play up the importance of the FDT more than I have. I
  dunno. -- Nov 11, 2020}.

\begin{mytheorem}[The First Derivative Test (Newtonian version)]
  Suppose $P(t)$ represents position as a function of time.
  \label{thm:FDT-Newton}
  \begin{enumerate}\label{list:FirstDerivTest-Newton}
  \item If the derivative of $P(t)$ is positive then the quantity being measured
    is, metaphorically, moving in the positive direction: It is
    increasing.
  \item If the derivative of $P(t)$ is negative then the quantity being measured
    is, metaphorically, moving in the negative direction; It is decreasing.
  \item If the derivative of $P(t)$ is neither positive nor negative -- that is,
    if it is zero -- then the object has stopped moving\aside{The
      metaphor actually breaks down a little bit here. It is hard to
      say that a vertically tossed ball is not moving at the top of
      its flight since it immediately begins to fall. Nevertheless, at
      the top of its flight its velocity is zero and it does ``stop''
      just for an instant.}.
  \end{enumerate}
\end{mytheorem}

The derivative of a curve, however we choose to
interpret it, can tell us a lot about the curve. If we use Newton's
kinematic view and pretend that the curve gives the distance from the
origin (position) of some object graphed against elapsed time then the
derivative gives us its velocity. Clearly then when the derivative is
positive the distance increases, and when the derivative is negative
the distance decreases.

We will revisit this theorem in various forms in the next chapter,
ending with its modern formulation. However we are not abandoning the
Newtonian view. We will use whatever tools allow us to think most
clearly about a given problem and sometimes this ``kinematic''
approach will be the one best suited to the problem at hand.

    %     Recall that in order to find out how high
    %     the tossed ball will go we first needed to discover how long it takes
    %     for the ball's velocity to drop to zero. This is, obviously, because
    %     the ball continues to rise until its velocity is zero. Thereafter it
    %     will fall. Once we know that the velocity is zero at time, say
    %     $t_{\text{max}},$ then we have $v(t_{\text{max}})=0,$ and the
    %     position, $p(t_{\text{max}}),$ at the same time is the maximimum
    %     height of the ball's flight.

    %     This is all fairly intuitive. 

    %     \begin{wrapfigure}[]{r}{2in}
    %     \captionsetup{labelformat=empty}
    %     \centerline{\includegraphics*[height=2in,width=2in]{../Figures/metaphor1}}
    %     \label{fig:metaphor1}
    %     \end{wrapfigure}
    %     The graph represents on of the author's (purely subjective sense of
    %     the) level of anxiety plotted against the number of chocolate bars he
    %     owns. That is, it graphs anxiety $(A)$ as a function of chocolate bars:
    %     $A(c).$

    %     Can we use this same kind of reasoning to find the coordinates of the
    %     highest point on the graph at the right?

    %     Notice that the horizontal and vertical axes are no longer
    %     ``position'' and ``time.'' Since there is no motion taking place there
    %     is clearly no velocity, so we can't simply set our velocity function
    %     equal to zero an solve for the time, $t,$ when velocity is zero like
    %     we did for the falling ball. Not only is there no velocity function,
    %     there is no time $(t)$ either!

    %     Or can we? Yes, it is true that there is no motion, and no time
    %     involved. But so what? Once we plot velocity against time we just have
    %     a curve. The only place where either velocity or time appear is in the
    %     labelling of the axes. 

    %     So let's just change the labels.

    %     That is, let's pretend -- for now -- that this \emph{is} a motion
    %     problem, that the vertical axis \emph{is} position, and that the
    %     horizontal axis \emph{is} time. In that case we can proceed as before:
    %     Find the velocity function (derivative), set it equal to zero to find
    %     the time (number of chocolate bars) when the velocity (derivative) is
    %     zero, put that back into the position (anxiety) function, $A(C),$ to
    %     find the maximum level of anxiety.

    %     This was Newton's idea.  Any such problem can be thought of as a
    %     problem of motion. In this case we think of our anxiety level as
    %     dependent on how many chocolate bars we have. If the number of
    %     chocolate bars changes then the anxiety level increases (or decreases)
    %     accordingly.

    %     But it is silly to rename our functions just to do the analysis and
    %     then change them back. Rather than explicitly pretending that motion
    %     is involved we keep that conceit in the background and adapt our
    %     notation and language, using motion as a metaphor.

    %     In particular we have seen that $\dfdx{p}{t}=v,$ (velocity is the
    %     derivative of position). Since velocity is by definition the rate of
    %     change of position ``with respect to time'' we adopt the convention,
    %     in this case, of interpreting the symbol $\dfdx{A}{C}$ as the rate of
    %     change of anxiety ``with respect to the number of chocolate bars we
    %     own.''  

    %     In this way the metaphor of motion is simply built into the language
    %     and notation we use, and we say things like ``As the number of
    %     chocolate bars increases our anxiety decreases, then increases, and
    %     then decreases again,'' just as if $C,$  (the number of chocolate
    %     bars) is constantly increasing in the same way that $t,$ (time) does.


    %     \begin{wrapfigure}[]{l}{2in}
    %     \captionsetup{labelformat=empty}
    %     \centerline{\includegraphics*[height=2in,width=2in]{../Figures/metaphor2}}
    %     \label{fig:metaphor2}
    %     \end{wrapfigure}
    %     It can be illuminating to look at the graph of both our function and
    %     its derivative at the same time. The sketch at the left shows the graph
    %     of $A(C)$ (in green) and $\dfdx{A}{C}$ (in orange) together on the
    %     same set of axes.

    %     Since $A(0) = 5$ when $C=0$ we seem to be fairly anxious. This is probably
    %     because we have no chocolate. As our number of chocolate bars
    %     increases\footnote{Notice how the language of motion is naturally
    %     co-opted for our purposes. We either own chocolate or we
    %     don't. Nothing is changing, but we pretend that as time goes on we
    %     are getting more chocolate.}
    %     our anxiety level is decreasing so the value of $\dfdx{A}{C}$ is
    %     negative (below the horizontal axis). When $\dfdx{A}{C} =0,$ at
    %     slightly less than one chocolate bar, our anxiety level ``bottoms
    %     out.'' Presumably this is because we're looking forward to eating our
    %     chocolate. As we gain chocolate our anxiety level rises again,
    %     probably because in these uncertain times you never know who is
    %     coveting your chocolate, or when it might be taken from you just because you
    %     have a lot of it. As we accumulate more our anxiety level drops again
    %     because we have plenty to share so we don't need to worry about some
    %     thief taking it from us. 

    %     This is a silly example of course, but the point of this silliness is
    %     quite serious: 

    %     On the other hand, if we take Leibniz's geometric viewpoint that the
    %     derivative of a curve at a point gives us the slope of the graph of
    %     some function at that point, then a positive derivative implies an
    %     increasing function and a negative derivative implies a decreasing
    %     function.

    %     Even the differential equation we solved in the previous chapter are
    %     amenable to using motion as a metaphor as the following example shows. 

    %     \myexample{Population Growth}
    %     \label{exam:population-growth}
    %     Suppose we have a single microbe living in a petri dish and that it
    %     divides once every hour. Then after one hour we have two microbes and
    %     our population has grown by one microbe per hour. After the second
    %     hour our population is now $4$ microbes, and the rate of growth during
    %     that second hour is $2$ microbes per hour. If this continues for $10$
    %     hours our microbe population will have grown according to the
    %     following chart\marginpar{This is supposed to be $P(t)=(5/2)^t.$ Check the numbers.}:
    %     $$
    %     \begin{array}{|ccc||ccc|}
    %     \hline
    %     \text{hours}&\text{population}&\text{rate of
                                          %                                           growth}&\text{hours}&\text{population}&\text{rate
                                                                                                                              %                                                                                                                               of
                                                                                                                              %                                                                                                                               growth}\\\hline{}
    %     1&2.5&1.5\frac{\text{\tiny microbes}}{\text{\tiny hour}}&6&244.1&97.2\frac{\text{\tiny microbes}}{\text{\tiny hour}}\\[1mm]
    %     2&6.3&2.5\frac{\text{\tiny microbes}}{\text{\tiny hour}}&7&610.4&242.9\frac{\text{\tiny microbes}}{\text{\tiny hour}}\\[1mm]
    %     3&15.6&6.2\frac{\text{\tiny microbes}}{\text{\tiny hour}}&8&1525.9&607.2\frac{\text{\tiny microbes}}{\text{\tiny hour}}\\[1mm]
    %     4&39.1&15.5\frac{\text{\tiny microbes}}{\text{\tiny hour}}&9&3814.7&1518.0\frac{\text{\tiny microbes}}{\text{\tiny hour}}\\[1mm]
    %     5&97.7&38.9\frac{\text{\tiny microbes}}{\text{\tiny hour}}&10&9536.7&3795.1\frac{\text{\tiny microbes}}{\text{\tiny hour}}\\[1mm]\hline
    %   \end{array}
    %     $$
    %     Can we find a function $P(t)$ which  tells us our microbe
    %     population at any time?

    %     \begin{wrapfigure}[]{l}{2in}
    %     \captionsetup{labelformat=empty}
    %     \centerline{\includegraphics*[height=1in,width=2in]{../Figures/velocity}}
    %     \label{fig:velocity}
    %     \end{wrapfigure}


    % %     \begin{wrapfigure}[]{R}{2in}
    %     \includegraphics*[height=1.75in,width=2in]{../Figures/parab-sec}
    % %     \caption{}
    % %     \label{fig:parab-sec}% \end{wrapfigure}

    %     \centerline{\includegraphics*[height=2in,width=2in]{../Figures/parab-tan}}

    % % %     \begin{wrapfigure}[]{L}{2in}
    %     \centerline{\includegraphics*[height=2in,width=2in]{../Figures/ProposedGraph}}
    % % %     \caption{\textcolor{blue}{Proposed graph of $h=h(t)$}}
    % % %     \label{fig:proposed}
    % % %   \end{wrapfigure}

    %     Because the change in velocity, $v,$ is so simple, $v=15-9.8t$ it is easy to
    %     discover \emph{when} the ball reaches it maximum height,
    %     but what \emph{is} the maximum height?  This is a bit trickier, as the
    %     change in position is not as simple.  Notice that
    %     initially the velocity was
    %     \(15 m/s.\) If the ball traveled that speed throughout the first
    %     second (which it doesn't), then the ball would travel \(15\)
    %     meters in the first second.  At one second, the velocity was \(5.02\,m/s,\)
    %     so if it stayed that velocity (which it doesn't), then in the
    %     next second the ball would again travel $5.02$) meters. But if the
    %     ball is accelerating it will actually go further during the second
    %     $1-$second interval than the first. We need to figure out some way to
    %     measure the distance traveled for an object whose velocity is
    %     continually changing.


    %     We will develop the tools for this soon.  But first let's look at this
    %     problem graphically.  If we plot the velocity
    %     $v=15-9.8t$ as a function of time, we have 

    %     Notice that the slope of this line is $-9.8,$ the (constant) acceleration due
    %     to gravity.  This makes sense since the slope of this line is the
    %     change in velocity divided by the change in time
    %     ($\frac{\text{m}}{\text{s}}/s = m/s^2$), which as we have seen is the
    %     acceleration, and Galileo showed that the acceleration is
    %     constant\footnote{Actually it isn't. The acceleration of falling
    %     objects depends on their distance from the center of the
    %     Earth. However, \emph{near} the surface of the Earth the change is
    %     so small that Galileo couldn't measure it with the tools he had. In
    %     any case, the change is so slight that it can be taken to be
    %     constant for our purposes.}.


    %     But what about the graph of the height, \(h(t)?\)
    %     A graph of what \(h=h(t)\)
    %     might look like is given in the following figure:

    %     \centerline{\includegraphics*[height=2in,width=2in]{../Figures/ProposedGraph}}

    %     This is not the trajectory of the ball (remember it is moving straight
    %     up and down), but is a plot of height versus time.  Since the velocity
    %     is continually changing, the graph of the height versus time will not have a constant slope the
    %     way that the graph of velocity did.

    %     There is a relationship between
    %     the graphs in these figures which is easiest to see by considering
    %     what the horizontal and vertical axes measure. In the graph
    %     \centerline{\includegraphics*[height=1in,width=2in]{../Figures/velocity}}
    %     the vertical axis is the velocity which is measured in meters per
    %     second, $m/s$ and the horizontal axis is measured in seconds.
    %     In the graph
    %     \centerline{\includegraphics*[height=2in,width=2in]{../Figures/ProposedGraph}}
    %     the vertical axis is the height, $h,$ which is measured in meters, and again
    %     the horizontal axis is the time, $t,$ which is measured in seconds.
    %     The first graph measures velocity directly. However Since the velocity is the change in
    %     distance divided by the change in time, we can use the information in
    %     the second graph to deduce the velocity at some time $t_0$ as
    %     follows.  We follow Leibniz's example and think of the graph as made
    %     up of infinitely small straight line segments, infinitesimals. In that
    %     case the slope of the second graph at $t_0$ is given by $\dfdx{h}{t}.$
    %     Since this fraction is measured in meters per second ($m/s$) the slope
    %     of the graph of height versus time is the velocity of the ball at the
    %     time $t_0.$ Or, to put it another way, the vertical coordinate of the
    %     first graph (velocity versus time) at any time $t_0$ gives the slope
    %     of the the second graph at the same time, $t_0.$

    %     More abstractly, we can say that the velocity (vertical coordinate of
    %     velocity versus time) at time $t_0$ is the slope of the line tangent
    %     to the graph of position (height versus time) at time $t_0.$
    %     This abstraction is what we have really been aiming for all along
    %     however before we go too far down the rabbit hole of abstraction,
    %     let's review what we've learned from this very concrete physical
    %     problem.


    %     then it seems that the slope
    %     of the curve in the second at any point \((t,h(t))\) should be the
    %     value \(v(t).\) But what do we mean by the slope of a curve?  This
    %     where our physical depiction of throwing a ball into the air might
    %     help.


    %     As we noted, on the way up the ball is continually slowing down.  On
    %     the way down, it is continually speeding up.  We can find the average
    %     velocity over any time interval \([t_0,t_1]\)
    %     by computing the ratio \(\frac{h(t_1 )-h(t_0 )}{t_1-t_0}.\)
    %     On our proposed graph of \(h(t),\) 
    %     this would be the slope of the line connecting the points
    %     \((t_0,h(t_0))\)
    %     and \((t_1,h(t_1))\)

    %     \centerline{\includegraphics*[height=1.75in,width=2in]{../Figures/parab-sec}}


    %     \begin{embeddedproblem}{}
    %     What would happen to the
    %     slope of the line segment joining \((t_0,h(t_0 ))\) to \((t_1,h(t_1 ))\) in
    %     relation to the slope of the tangent line?  How does this translate
    %     into the velocities of the ball in the original problem?
    %     \end{embeddedproblem}

    %     However, as we noted before, this is not the velocity of the ball at
    %     time \(t_0\) since that ball is slowing down.

    %     If we could somehow force the ball to stop slowing down, if we could
    %     somehow turn off gravity, at the instant \(t_0\) then from that time
    %     on the ball would move in a straight line with the velocity it has at
    %     the instant \(t_0.\) A moment's reflection makes it clear that this
    %     line must be the tangent line to the curve \(h(t)\) at the point
    %     \((t_0,h(t_0 )).\)


    %     Figure 5

    %     Bud – later when we talk about limits, we can come back to this and use a series of drawings where the time interval gets smaller.  Or do you think we should get into this now?  One possibility is the following problem

    %     While it is true that we can't physically make the ball stop
    %     slowing down, this 'thought experiment' gives us insight into
    %     the relationship between the graph of the ball's velocity and height
    %     versus time. Clearly the
    %     values of \(v(t)\) represent the slopes of the tangent lines to the curve
    %     \(h=h(t).\)

    %     \begin{wrapfigure}[]{L}{2in}
    %     \includegraphics*[height=2in,width=2in]{../Figures/parab-tan}
    %     \caption{}
    %     \label{fig:parab-tan}
    %     \end{wrapfigure}



    %     Using the tools of  Calculus we can exploit this relationship between
    %     the graph of a curve and the slope of its tangent line.  This
    %     fact led to one of ``golden periods'' of mathematics in the
    %     1700s where the power of Calculus was used to examine all sorts of
    %     geometric and physical problems.  In order to exploit this
    %     relationship, we will need to establish rules for computing.  Indeed,
    %     as we observed in the previous section, this is what a 'calculus' is a
    %     set of rules for computation.  In the case of differential Calculus,
    %     it is literally a set of rules for computing infinitely small
    %     (infinitesimal) differences.  The key turns out to be relating
    %     everything back to the graphs of straight lines since straight lines
    %     are the graphs of quantities whose rates of change are constant and
    %     thus are easily measured.

    %     Developing these rules will allow us to examine the relationship
    %     between the graphs of \(v=v(t)\)
    %     and \(h=h(t)\)
    %     algebraically.  In particular, we will see how to find the equation of
    %     one curve, for example the height of a falling object, if we are given
    %     the other, for example its velocity.  Once we develop these tools, we
    %     will be ready to see how Calculus led to this golden period of
    %     mathematics. 

    %     \subsection{Tangent Lines}
    %     \label{sec:tangent-lines}

    %     In many ways, the invention of calculus was the culmination of the
    %     merging of geometry and kinematics with algebra, while exploiting
    %     infinitesimals.  While it is true that Algebraic methods are the
    %     predominant tool of a calculus class, the blending of it with geometry
    %     and kinematics helps students adapt these techniques to physical and
    %     geometric applications.  In a curriculum that has downplayed the role
    %     of geometry, having precalculus problems which bring it back can pay
    %     future dividends.

    %     Consider the method of Gilles Personne de Roberval (1602-1675) for
    %     constructing the tangent to a curve.  In a blend of kinematics and
    %     geometry, Roberval considered a curve as generated by a point whose
    %     motion is compounded from 2 known motions.  The ``velocity vectors'' of
    %     the motions will combine to give the resultant tangent vector to the
    %     curve [Boyer, p. 390].  For example, consider that a parabola is the
    %     set of points in the plane equidistant from a fixed point (focus) and
    %     fixed line (directrix).  We can think of the parabola as generated by
    %     point moving so that it maintains the same distance from the focus as
    %     it does from the directrix.\\
    %     \centerline{\includegraphics*[height=3in,width=2.5in]{../Figures/Roberval1}}

    %     Since the point maintains the same distance from the focus as it does
    %     from the directrix, the velocity vectors of these two motions (away
    %     from $F$ and $d$) will be the same length and the tangent vector will be
    %     the diagonal of the rhombus.  While this method is limited from an
    %     Algebraic point of view (in terms of computing a slope), it does have
    %     a nice geometric consequence.  Since the diagonal of a rhombus bisects
    %     the angle, then we have $\angle 1\cong\angle2\cong\angle3.$  This provides the reflective
    %     property of a parabola which says that any light ray, sound wave,
    %     radio wave, etc. parallel to the axis of the parabola will be
    %     reflected to the focus.  This is important to the design of satellite
    %     dishes and reflective telescopes. \\
    %     \centerline{\includegraphics*[height=3in,width=5in]{../Figures/Telescopes}}
    % \begin{verbatim}
    %     First image from http://media-2.web.britannica.com/eb-media/82/72982-004-9191583E.jpg. 
    %     Second image from http://astro-canada.ca/_en/_illustrations/a4304_newton_en_p.jpg.
    % \end{verbatim}

    %     \begin{embeddedproblem}{}
    %     Given that an ellipse is the set of points in the plane, the sum of
    %     whose distances from two fixed points (foci) is constant, use
    %     Roberval's idea to construct the tangent to the ellipse.  Notice
    %     that one of the motions is away from a focus and the other is
    %     towards the other focus.  Why?\\
    %     \centerline{\includegraphics*[height=2in,width=2in]{../Figures/Roberval2}}

    %     Conclude that the parallelogram must be a rhombus and use this to
    %     demonstrate the reflective property of the ellipse, namely that any
    %     ray emanating from one focus of the ellipse is reflected by the
    %     ellipse to the other focus.   
    %     \end{embeddedproblem}

    %     The aforementioned reflective property is noticed in a whispering
    %     gallery, where a person at one focus can speak quietly and the person
    %     at the other focus can hear it.  It is also used in the treatment of
    %     kidney stones through a procedure called extracorporeal shock wave
    %     lithotripsy (ESWL).  As imposing as the title sounds, the name makes
    %     sense.  Lithotripsy literally means ``rock grinding'' and the procedure
    %     utilizes shockwaves generated outside the body to crush kidney stones
    %     into granules that can pass relatively painlessly.  Ellipses are
    %     utilized to direct the shock waves from the generator to the stone as
    %     can be seen below.\\
    %     \centerline{\includegraphics*[height=2in,width=2in]{../Figures/Lithotripsy}}

    %     Image 1 from http://www.hussainkidney.com/images/LithotripsyVSsurgery.jpg.
    %     Image 2 from http://blogs.nejm.org/now/wp-content/uploads/2012/07/Electrohydraulic-Lithotripter.jpg

    %     As can be seen from the images, shock waves are generated outside the
    %     body at one focus of the elliptical reflector and are reflected to the
    %     other focus where the stone is located.  The procedure is done on an
    %     outpatient basis and is a very effective, non-invasive procedure for
    %     treating kidney stones.

    %     As another (future) application, consider the following image.
    %     \centerline{\includegraphics*[height=2in,width=2in]{../Figures/MagRes}}

    %     Image from http://physics.bu.edu/~condmat/pictures/mirage.jpg .

    %     This image represents the magnetic resonance of 36 cobalt atoms placed
    %     in an elliptical ring called a quantum corral by scientists at IBM
    %     Almaden labs.  When researchers bombarded a cobalt atom placed at one
    %     focus with electrons, the electron partial waves reflected off of the
    %     elliptical ring to produce a weaker ``mirage'' at the other focus even
    %     though there was no atom present
    %     [http://cerncourier.com/cws/article/cern/28197].  IBM is hoping to
    %     utilize this ``Quantum Mirage'' effect to transmit information on a
    %     nanoscale, where normal circuitry does not work.  This would lead
    %     toward nano-sized processors.  Below is an artist's rendition of a
    %     nanobot attaching itself to a damaged nerve and becoming a prosthetic
    %     nerve
    %     [http://www.cg4tv.com/animation/stock-images-3d/neurons-nanobot-high-resolution.jpg].
    %     \\
    %     \centerline{\includegraphics*[height=2in,width=2in]{../Figures/Nanobot}}
    %     This technology doesn't yet exist, but it demonstrates how even old
    %     ideas can lead to future developments.

    %     \begin{embeddedproblem}{}
    %     A hyperbola is the set of points in the plane, the difference of
    %     whose distances from two fixed points (foci) is constant.  Use
    %     Roberval's idea to construct the tangent to the hyperbola and derive
    %     the reflective property of a hyperbola, namely that any ray directed
    %     at one focus will reflect off of the hyperbola to the other focus.\\
    %     \centerline{\includegraphics*[height=2in,width=2in]{../Figures/Roberval3}} 
    %     \end{embeddedproblem}

    %     The aforementioned reflective property is utilized in the secondary
    %     reflector of a cassegrain antenna.\\
    %     \centerline{\includegraphics*[height=2in,width=2in]{../Figures/Telescope2}}
    %     Image 1 from http://www.antedo.com/images/13-Meter-(TCS).jpg .
    %     Image 2 from http://www.rfwireless-world.com/images/cassegrain-feed-antenna.jpg .


    %     \subsection{}

    %     \begin{wrapfigure}[]{O}{3in}
    % %       \vskip-.7cm{}
    %     \includegraphics*[height=1.75in,width=3in]{../Figures/CosApprox2}
    %     \caption{\textcolor{blue}{Approximating $\cos x$ with a quadratic polynomial}}
    %     \label{fig:CosApprox2}
    %     \end{wrapfigure}



    %     \begin{wrapfigure}[]{O}{3in}
    % %       \vskip-.7cm{}
    %     \includegraphics*[height=1.75in,width=3in]{../Figures/CosApprox4}
    %     \caption{\textcolor{blue}{Approximating $\cos x$ with a quartic polynomial}}
    %     \label{fig:CosApprox4}
    %     \end{wrapfigure}

    %     \marginpar{These last paragraphs should be somewhere else. But where?}
    %     For example, for values of \(x\)
    %     between \(-\frac\pi2\) and \(\frac\pi2\) the
    %     \label{page:cosine-approx}
    %     cosine of \(x\) is closely approximated by the polynomial
    %     \(1-\frac{x^2}{2}+\frac{x^4}{24}.\) So for values of \(x\) in that
    %     range we can replace \(\cos(x)\) with
    %     \(1-\frac{x^2}{2}+\frac{x^4}{24}\) with very little loss of
    %     accuracy. This is useful because, for example, \(\cos(.7)\) can be
    %     very hard to compute whereas \(1-\frac{(.7)^2}{2}+\frac{(.7)^4}{24}
    %     \approx 0.77\) is relatively easy.

    %   %     Once known this approximation is easy to confirm and to use.  But
    %   %     coming up with it in the first place, finding the right polynomial
    %   %     to use is quite difficult, until you've learned a little Calculus as
    %   %     we will see in section~\ref{subsec:tangent-lines}.

    %   %     This approximation is easy to confirm and to use \emph{once it is
    %   %     known.}  But coming up with it in the first place is quite
    %   %     difficult until you've learned a little Calculus as we will see in
    %   %     section~\ref{subsec:tangent-lines}.\marginpar{A word about the use
    %   %     of technology should be added here.}

    %     \subsection{}
    %     Of particular note is the following technique of
    %     Gilles Personne de Roberval (1602-1675).  Roberval thought of a
    %     curve in the plane as being generated by a point whose motion is
    %     compounded from two known motions.  By combining the ``velocity
    %     vectors'' of these motions, the tangent vector will
    %     result. With the tangent vector in hand the tangent line is easily
    %     computed.


    %     We'll apply this idea to a parabola.  You probably think of a parabola
    %     as being determined by a quadratic function (which it is), but that is
    %     the modern approach. The older way to define a parabola is geometric: A
    %     parabola is the set of points in a plane which are equidistant from a
    %     fixed point called the focus and a fixed line called the directrix. In
    %     the figure, the point $P$ on the parabola is the same
    %     distance from the focus $F$ as it is from the directrix $d.$ These
    %     lengths  change as $P$ moves along the parabola, but the distance
    %     from $P$ to $F$ will always equal the distance from $P$ to $d.$
    %     \begin{embeddedproblem}{}
    %     \begin{description}
    %     \item[a.] Explain why the parallelogram in the figure is
    %     a rhombus.
    %     \item[b.] Given that the parallelogram in the  figure is
    %     a rhombus, show that the two angles $\alpha,$ and $\beta$ are
    %     congruent.
    %     \end{description}
    %     \end{embeddedproblem}

    %     If we think of the motion of $P$ as a combination of motion away from
    %     $F$ and motion away from $d,$ then we can modify our drawing to
    %     reflect this as in the figure. While this geometric approach does
    %     not produce the equation of the tangent line, it does have the
    %     interesting consequence
    %     \begin{wrapfigure}[]{O}{2in}
    % %       \vskip-.7cm{}
    %     \includegraphics*[height=1.5in,width=2in]{../Figures/Ellipse}
    %     \caption{}
    %     \label{fig:ellipse}
    %     \end{wrapfigure}
    %     that any light ray, radio wave, sound
    %     wave, etc. which comes in parallel to the axis of the parabola
    %     (perpendicular to the directrix) will reflect off of the parabola to
    %     the focus.  Among other things, this is why satellite dishes and
    %     primary mirrors in reflective telescopes are parabolic.  
    %     \begin{embeddedproblem}{}
    %     Consider that an ellipse is the set of points in a plane, the sum of
    %     whose distances from two fixed points (foci) is
    %     constant. Using Roberval's idea, we can think of the
    %     motion of a point P on the ellipse as being the result of a motion
    %     away from one focus and toward the other focus.
    %     Explain why the parallelogram in the figure must be a rhombus and
    %     use this to show that angle $\alpha$ is congruent to angle $\beta.$
    %     \end{embeddedproblem}

    %     \begin{wrapfigure}[]{O}{2in}
    % %       \vskip-.7cm{}
    %     \includegraphics*[height=1.5in,width=2in]{../Figures/RockGrinding}
    %     \caption{}
    %     \label{fig:RockGrind}
    %     \end{wrapfigure}

    %     This says that any light ray, sound wave, etc. emanating from
    %     one focus will reflect off of the ellipse to the other focus.  Among
    %     other things, this reflective property is used to treat kidney stones.
    %     The treatment is called Extracorporeal Shock Wave Lithotripsy (ESWL)
    %     and is really what the name says.  Lithotripsy literally means ``rock
    %     grinding,'' and this technique uses shock waves which are generated
    %     outside the body (See
    %     Figures~\ref{fig:RockGrind})\footnote{http://us.medispec.com/patient/kidney-stones/treatment-using-eswl/
    %     and\\
    %     http://bme240.eng.uci.edu/students/09s/ysantoro/CurrTechn.html} As
    %     you can see in the figure, the table contains a cup which is in
    %     the shape of part of an ellipse.  At one focus of the ellipse is an
    %     electrode which generates shock waves.  These shock waves reflect off
    %     of the elliptical cup and converge on the other focus.  When the
    %     kidney stone is positioned at that focus, the shock waves will pummel
    %     it into small pieces which can be passed relatively painlessly through
    %     the ureter.
    %     \begin{embeddedproblem}{}
    %     An hyperbola is the set of points in the plane, the difference of
    %     whose distances from two fixed points (foci) is constant.  Using the
    %     ideas of Roberval, show that any light ray, sound wave,
    %     etc. directed at one focus will reflect off of the hyperbola toward
    %     the other focus.
    %     \end{embeddedproblem}
    %     This idea is utilized in the secondary mirror of Cassegrain
    %     telescopes.

    %     \begin{wrapfigure}[]{O}{2in}
    % %       \vskip-.7cm{}
    %     \includegraphics*[height=2in,width=2in]{../Figures/ReflectingHyperbola}
    %     \caption{}
    %     \label{fig:ReflectingHyperbola}
    %     \end{wrapfigure}

    %     Roberval's idea to use distances from fixed objects to describe the
    %     location of a point on a plane is the essence of analytic geometry.
    %     You are accustomed to plotting points by determining the distances
    %     from two fixed axes to give the coordinates of a point.  This allows
    %     us to identify geometric curves by their corresponding equations.
    %     Specifically, a point will lie on a curve provided its coordinates
    %     satisfy a certain equation.  For example, the point $P$ with coordinates
    %     $(2,4)$ lies on the parabola with equation $y=x^2$ since $4=2^2,$ whereas
    %     the point $Q$ with coordinates $(2,5)$ does not since $5\neq2^2.$
    %     Notice that, remarkably, we were able to answer a geometry question
    %     without drawing a picture!  We're not promoting not drawing pictures,
    %     but it does illustrate the power of analytic geometry; we can bring
    %     the tools and techniques of Algebra to bear on solving the geometry
    %     problems.  

    %     \begin{embeddedproblem}{}
    %     Before the invention of calculus, Galileo tried to determine the
    %     motion of a projectile as it fell to the ground.  Efforts of people
    %     such as Galileo in handling specific problems such as this helped
    %     Newton and Leibniz to generalize the ideas to form differential
    %     calculus.  We will recreate some of these ``pre-calculus'' ideas.
    %     To start, by experimenting with a ball rolling down a plane, Galileo
    %     surmised that an objected falling under the influence of gravity
    %     accelerated at a constant rate (ignoring air resistance, friction,
    %     etc.).  Today we know that rate to be $9.8$ (meters/second)/second.
    %     Suppose we stand at some height and drop a ball.  This says that the
    %     velocity of the object at time t seconds is given by $v(t)=9.8t$
    %     meters/second.  Figuring out how far the ball has fallen is a bit
    %     trickier as the velocity is not constant.
    %     \begin{description}
    %     \item[(a)] Suppose we divide the time interval $[0,t]$ into four
    %     equal intervals.  Determine that average value of $v(0), v(t/4),
    %     v(2t/4), v(3t/4), v(4t/4)$ is $4.9t.$ What if we divided the time
    %     interval into $8, 16,$ or $32$ equal intervals?
    %     \item[(b)] Assuming that the average velocity during the time
    %     interval $[0,t]$ is given by $4.9t\,m/s,$  show that the distance
    %     the ball has fallen at time $t$ is given by $y(t)=4.9t^2.$
    %     \item[(c)] To check if Galileo is correct that $y(t)=4.9t^2$ meters,
    %     compute the average velocity of the ball from time $t$ to time
    %     $t+\Delta t$ for $\Delta t=.1, .01, .001,$ and $.0001.$
    %     \item[(d)] Based on what you obtained in part (c), what would the
    %     instantaneous velocity be at time $t,$ (i.e., when $\Delta t=0$).
    %     Does this agree with what we had as $v(t)$ before?
    %     \end{description}
    %     \end{embeddedproblem}


    %     \begin{myexample}
    %     Consider the curve $y=x^2:$\\
    %     \centerline{\includegraphics*[height=2in,width=2in]{../Figures/Quadratic}}
    %     We wish to find an equation of the line tangent to this curve at
    %     the point $(3,9).$
    %     It is clear from the graph that the line tangent to this curve at
    %     the point $(3,9)$ will also pass through the $y-$axis, say at
    %     the point $(0,b).$ If we can find $b$ we have the tangent line. 

    %     The slope of the tangent line is $m=\frac{9-b}{3}$ so the equation
    %     of the tangent line will be 
    %     \begin{equation}
    %     y=\left(\frac{9-b}{3}\right)x+b\label{eq:TanSq}
    %     \end{equation}
    %     for one, and only one, value of $b.$ If $b$ is anything other than
    %     that one value then equation(~\ref{eq:TanSq}) will intersect the parabola
    %     twice: once at the point $(3,9),$ and once at another point, say
    %     $(x,x^2).$ At this second point we have 
    %     $$
    %     x^2=\left(\frac{9-b}{3}\right)x+b.
    %     $$

    %     Rearranging this slightly we have 
    %     $$
    %     x^2-\left(\frac{9-b}{3}\right)x-b=0
    %     $$
    %     which is a quadratic and we can solve quadratics.

    %     Before charging blindly forward though, let's stop and think about
    %     what we're trying to accomplish. We want to find the value of $b$
    %     which guarantees that there is only \emph{one} solution. We're
    %     not really interested in the actual solution itself. We just want
    %     to ensure that there is only one.

    %     From the Quadratic Formula we see that:
    %     $$
    %     x=
    %     \frac{\left(\frac{9-b}{3}\right)\pm\sqrt{\left(-\left(\frac{9-b}{3}\right)\right)^2+4b}}{2},
    %     $$
    %     which has one solution precisely when the discriminant\footnote{The part under
    %     the square root symbol.} is zero. Thus
    %     \begin{align*}
    %     \left(\frac{9-b}{3}\right)^2+4b&=0\\
    %     \frac{9^2-18b+b^2}{9}+4b&=0\\
    %     9^2-18b+b^2+36b&=0\\
    %     b^2+18b+9^2&=0\\
    %     (b+9)^2&=0,
                   %   \end{align*}
                   %                    so $b=-9.$ Thus an equation of the line tangent to $y=x^2$ at the point
                   %                    $(3,9)$ is
                   %                    $$
                   %                    y=6x-9.
                   %                    $$



                   %                    \begin{embeddedproblem}{}
                   %                    Use this method to find an equation of the line tangent to
                   %                    $y=x^2$ at each of the following points:\\
    %     \begin{description}
    %     \item[(a)] $(2,4)$
    %     \item[(b)] $(-3,9)$
    %     \item[(c)] $(1,1)$
    %     \item[(d)] $\left(\frac12,\frac14\right)$
    %     \item[(e)] $(a,a^2)$
    %     \end{description}
    %     \end{embeddedproblem}  

    %     \begin{embeddedproblem}{}
    %     Use this method to find an equation of the line tangent to the given
    %     curve  at the given point:\\
    %     \begin{description}
    %     \item[(a)] $y=x^2+x$ at $(2,6).$
    %     \item[(b)] $y=x^2+x$ at $(\tau,\tau^2+\tau).$
    %     \item[(c)] $y=x^2-3x+2$ at $(\tau,\tau^2-3\tau+2).$
    %     \item[(d)] $y=ax^2+bx+c$ at $(\tau,a\tau^2+b\tau+c).$
    %     \end{description}
    %     \end{embeddedproblem}
    %   %     For example, he showed that the line tangent to the parabola at the
    %   %     point $(0,b)$ in the following diagram will always pass through the
    %   %     point $(0,-b).$
    %     \end{myexample}




%%% Local Variables: 
%%% mode: latex
%%% outline-minor-mode: t
%%% eval:(flyspell-mode)
%%% TeX-master: "DiffCalc22"
%%% End: 
